% Aporte de Jilmar Cedeño Intriago. Incluir desde el informe principal.
\subsection{Implementación con sockets TCP}
\label{sec:sockets-tcp}

Usamos sockets TCP porque permiten intercambiar un flujo de bytes fiable y ordenado~\cite{eddy2022tcp}. Codificamos cada mensaje como JSON UTF-8 y anteponemos una cabecera de cuatro bytes con la longitud del cuerpo en orden de red. Para enviarlo usamos \texttt{sendall()} y, al recibirlo, leemos en bucle hasta completar la cabecera y el cuerpo. Así recuperamos mensajes completos aunque una lectura devuelva solo parte del contenido y podemos reutilizar la conexión.

Para la prueba en loopback levantamos tres nodos lógicos en un mismo proceso, cada uno con su puerto TCP: \texttt{nodo-A:5101}, \texttt{nodo-B:5102} y \texttt{nodo-C:5103}. Distribuimos las solicitudes en el ciclo $A\to B$, $B\to C$ y $C\to A$. Configuramos cada servidor para devolver la misma carga por la conexión persistente. En modo concurrente usamos una barrera para sincronizar el inicio de cada grupo de tres solicitudes.

Inicializamos en cero el contador lógico de cada nodo. Antes de enviar aplicamos $L = L + 1$. Cuando recibimos un mensaje con sello $T$, aplicamos $L = \max(L,T) + 1$. Protegemos ambas actualizaciones y el registro en memoria con el mismo bloqueo para que el cliente y el servidor del nodo compartan una secuencia lógica. Por cada intercambio completo registramos cuatro eventos: el envío y la recepción de la solicitud y de la respuesta. Guardamos también la hora UTC con precisión de microsegundos en un campo independiente.

Medimos el RTT en el cliente con \texttt{time.perf\_counter()} desde antes del envío hasta que decodificamos la respuesta completa. Incluimos así el procesamiento remoto y la serialización y decodificación del intercambio. Abrimos las conexiones antes de medir y exportamos los CSV al finalizar. Para calcular el rendimiento dividimos el número de intercambios completados por el tiempo total de la tanda, medido con el mismo reloj.

\subsubsection{Desarrollo, tipado y manejo de errores}

En nuestra implementación TCP controlamos el formato del mensaje y el flujo de lectura y escritura con la biblioteca estándar. Definimos el encuadre, validamos los tipos de los campos JSON y comprobamos que cada respuesta corresponda a la solicitud. También configuramos tiempos de espera para las operaciones del socket. Al comparar con gRPC, consideramos que este usa HTTP/2~\cite{thomson2022http2} y ofrece un contrato con Protocol Buffers, interfaces generadas y estados RPC. Tomamos el estudio de Niswar y colaboradores como antecedente de evaluación en microservicios~\cite{niswar2024microservices}.
